\documentclass[11pt]{amsart}
\IfFileExists{lmodern.sty}%
  {\usepackage[T1]{fontenc}\usepackage{lmodern}\usepackage{microtype}}%
  {\usepackage[expansion=false]{microtype}}
\usepackage{amssymb,mathtools}
\usepackage{booktabs,array,longtable}
\usepackage{enumitem}
\emergencystretch=2.5em
\DeclareMathSizes{10}{10}{7}{6}
\DeclareMathSizes{9}{9}{6}{6}
\DeclareMathSizes{8}{8}{6}{6}
\usepackage[colorlinks=true,linkcolor=blue,citecolor=blue,urlcolor=blue]{hyperref}
\hypersetup{pdftitle={Irrationality measures for dilogarithms of rational numbers},pdfauthor={Sean Litsey}}

% ---------------------------------------------------------------------------
% Location of the Lean formalization.
\newcommand{\FormalizationURL}{https://github.com/slitsey/catalan-irrationality-measure}
\newcommand{\FormalizationLink}{\url{\FormalizationURL}}
% ---------------------------------------------------------------------------

\newtheorem{theorem}{Theorem}[section]
\newtheorem{lemma}[theorem]{Lemma}
\newtheorem{proposition}[theorem]{Proposition}
\newtheorem{corollary}[theorem]{Corollary}
\theoremstyle{definition}
\newtheorem{definition}[theorem]{Definition}
\theoremstyle{remark}
\newtheorem{remark}[theorem]{Remark}
\numberwithin{equation}{section}
\setcounter{tocdepth}{1}

\newcommand{\Z}{\mathbb Z}
\newcommand{\Q}{\mathbb Q}
\newcommand{\R}{\mathbb R}
\newcommand{\C}{\mathbb C}
\newcommand{\F}{\mathbb F}
\newcommand{\N}{\mathbb N}
\newcommand{\Li}{\operatorname{Li}}
\newcommand{\lean}[1]{\nolinkurl{#1}}
\newcommand{\vG}{v_p^{\mathrm G}}
\newcommand{\Dlow}{\underline{D}}
\newcommand{\cont}{\operatorname{cont}}

\begin{document}

\title[Irrationality measures for $\Li_2(a/b)$]{Irrationality measures for dilogarithms of rational numbers}

\author{Sean Litsey}
\email{sean.litsey.math@gmail.com}

\date{November 18, 2026}

\subjclass[2020]{Primary 11J82; Secondary 11J72, 33B30, 15A15, 31A15, 68V20}

\keywords{Irrationality exponent, dilogarithm, Hankel determinants, $p$-adic valuations, logarithmic energy,
equilibrium measures, formal verification, Lean}

\begin{abstract}
Let $\mu$ denote the irrationality exponent. For coprime integers $a$ and $b$ with $0<2|a|\le b$ we prove an
explicit upper bound for $\mu(\Li_2(a/b))$ whenever an explicit margin $\eta(a,b)$ is positive. The margin is the sum of
a closed-form arithmetic rate in $\log b$, $\log|a|$ and $\log|b-a|$ and a piecewise affine decay rate in
$\log(b/|a|)$. For example $\mu(\Li_2(1/2))\le234.57$, $\mu(\Li_2(1/3))\le21.57$, $\mu(\Li_2(1/4))\le15.19$,
$\mu(\Li_2(1/5))\le12.92$, $\mu(\Li_2(-1/2))\le13.94$ and $\mu(\Li_2(-1/3))\le11.91$. These are, to our knowledge, the first irrationality measures for those six numbers.
For reduced numerator $2$ we obtain finite measures at $2/b$ for odd $b\ge13$ and at $-2/b$ for odd $b\ge11$.
The values at $-2/11$, $-2/13$ and $-2/15$ appear to give new irrationality results. For $6\le b\le20$ the bounds improve those of Rhin and Viola, for instance
from $783.29$ to $11.73$ at $b=6$ and from $11.94$ to $9.36$ at $b=20$. For $\Li_2(-1/b)$
with $5\le b\le20$ they improve the bounds that follow from Rhin and Viola's linear independence measures. The approximations are Tang's normalized Hankel determinants. Liu and Tan and Zhang
used them, with the same arithmetic rate, to prove irrationality of $\Li_2(\pm1/k)$ and of $\Li_2(p/q)$ in explicit
ranges. A measure needs more. We prove that the determinants do not vanish at rational points for a band of primes
of positive density, and we prove decay on an exponentially small neighbourhood of $\Li_2(a/b)$ by a dual
certificate for a logarithmic energy problem with capped lattice charges. Concavity of the dual and a transport in
$\log(b/|a|)$ make the decay bound uniform from finitely many certificates. The proof is formalized in Lean~4.
\end{abstract}

\maketitle
\tableofcontents

%% =====================================================================
\section{Introduction}\label{sec:intro}

\subsection{The result}
The \emph{irrationality exponent} $\mu(\xi)$ of a real number $\xi$ is the infimum of the real numbers $\mu$ such that
$|\xi-p/q|>q^{-\mu}$ for all integers $p$ and all sufficiently large integers $q$. It does not change when $\xi$ is
multiplied by a nonzero rational number or translated by a rational number. A finite bound for $\mu(\xi)$ is an
irrationality measure for $\xi$, and it shows that $\xi$ is irrational. The dilogarithm is
\[
 \Li_2(x)=\sum_{k\ge1}\frac{x^k}{k^2}\qquad(|x|\le1).
\]

Let $a$ and $b$ be coprime integers with $b>0$ and $0<2|a|\le b$. Put
\[
 \lambda=\frac ab,\qquad t=\log\frac b{|a|},\qquad s=\operatorname{sign}a,
\]
and define the \emph{arithmetic rate}
\begin{equation}\label{eq:crate}
 c(a,b)=\frac72-\frac{523}{840}+\frac83\log b-6\log|a|-\frac94\log|b-a| .
\end{equation}
Let $M=98/39$. In Section~\ref{sec:uniform} we define, from a finite table of certified data, an explicit function
$\Dlow_s(t,\rho)$ of the sign $s$, of $t\ge\log2$ and of $\rho>0$. It is the largest of finitely many numbers, each
of which is affine in $t$ on an interval. The table has $37$ nodes and $419$ entries. A summary is in
Appendix~\ref{app:table}, and the full table is part of the formal statement.

\begin{theorem}\label{thm:main}
Let $a,b$ be coprime integers with $b>0$ and $0<2|a|\le b$, and let $\rho>0$. Suppose that
$\eta=\Dlow_s(t,\rho)+c(a,b)>0$. Then for every
\[
 \mu>\rho\cdot\max\Bigl(\frac2\eta,\frac1M\Bigr)
\]
there is an integer $q_0$ such that
\[
 \Bigl|\Li_2\Bigl(\frac ab\Bigr)-\frac pq\Bigr|>q^{-\mu}
\]
for all integers $p$ and all integers $q\ge q_0$. In particular $\mu(\Li_2(a/b))\le\rho\max(2/\eta,1/M)$.
\end{theorem}

Theorem~\ref{thm:main} has been verified in the Lean~4 proof assistant, together with every theorem, proposition and
lemma used in its proof. The formal proof has no unproved steps and uses only the standard axioms.
Section~\ref{sec:lean} describes it. The remarks are not part of the formal proof. The use of AI in this work is
disclosed at the end of the paper.

The numbers in the following corollary are evaluations of the formula of Theorem~\ref{thm:main} at one admissible
$\rho$. They are rounded up in the last digit.

\begin{corollary}\label{cor:table}
For the numbers $x$ in Tables~\ref{tab:values} and~\ref{tab:larger-numerators}, and for $x=\Li_2(\pm2/39)$,
the stated bound $\mu(x)\le\mu_0$ holds in the plain form of Theorem~\ref{thm:main}. In particular, $\Li_2(2/b)$
has a finite irrationality measure for every odd $b\ge13$, and $\Li_2(-2/b)$ for every odd $b\ge11$.
\end{corollary}

\begin{table}[ht]
\centering
\renewcommand{\arraystretch}{1.1}
\setlength{\tabcolsep}{4pt}
\begin{tabular}{@{}r|rrr|rrr@{}}
\toprule
 & \multicolumn{3}{c|}{$x=\Li_2(1/k)$} & \multicolumn{3}{c}{$x=\Li_2(-1/k)$}\\
$k$ & $\mu_0$ & RV & Hata & $\mu_0$ & RV & Hata\\
\midrule
$2$ & $234.57$ & -- & -- & $13.94$ & -- & --\\
$3$ & $21.57$ & -- & -- & $11.91$ & -- & --\\
$4$ & $15.19$ & -- & -- & $10.99$ & -- & --\\
$5$ & $12.92$ & -- & -- & $10.38$ & $62.69$ & $228.61$\\
$6$ & $11.73$ & $783.29$ & -- & $10.02$ & $44.51$ & $55.94$\\
$7$ & $11.03$ & $69.69$ & $95.06$ & $9.76$ & $33.76$ & $35.84$\\
$8$ & $10.50$ & $47.43$ & $51.08$ & $9.58$ & $26.06$ & $27.39$\\
$9$ & $10.15$ & $30.44$ & $33.13$ & $9.43$ & $21.88$ & $22.66$\\
$10$ & $9.86$ & $24.04$ & $25.24$ & $9.31$ & $19.24$ & $19.75$\\
$11$ & $9.67$ & $19.80$ & $21.19$ & $9.21$ & $17.42$ & $17.77$\\
$12$ & $9.54$ & $17.60$ & $18.54$ & $9.12$ & $16.08$ & $16.33$\\
$13$ & $9.41$ & $16.29$ & $16.65$ & $9.16$ & $15.05$ & $15.24$\\
$14$ & $9.32$ & $15.13$ & $15.27$ & $9.14$ & $14.23$ & $14.42$\\
$15$ & $9.32$ & $14.40$ & $14.48$ & $9.16$ & $13.58$ & $13.75$\\
$16$ & $9.33$ & $13.71$ & $13.81$ & $9.17$ & $13.06$ & $13.20$\\
$17$ & $9.34$ & $13.18$ & $13.24$ & $9.21$ & $12.61$ & --\\
$18$ & $9.32$ & $12.71$ & $12.76$ & $9.23$ & $12.23$ & --\\
$19$ & $9.34$ & $12.30$ & -- & $9.27$ & $11.90$ & --\\
$20$ & $9.36$ & $11.94$ & -- & $9.30$ & $11.62$ & --\\
\bottomrule
\end{tabular}
\caption{Bounds $\mu(x)\le\mu_0$ from Theorem~\ref{thm:main}, with the published bounds. For $\Li_2(1/k)$ the
column RV is Rhin and Viola \cite[pp.~434--435]{RhinViola05}. For $\Li_2(-1/k)$ it is $\nu+1$, where $\nu$ is their
linear independence measure of $1$, $\Li_1(-1/k)$, $\Li_2(-1/k)$ \cite[Table~1]{RhinViola19}. The columns Hata are
\cite[Tables~1 and~2, p.~386]{Hata93dilog}. Published values are rounded to two decimals. Theorem~\ref{thm:main}
also gives $\mu(\Li_2(2/39))\le23.11$ and $\mu(\Li_2(-2/39))\le22.31$.}\label{tab:values}
\end{table}

\begin{table}[ht]
\centering
\begin{tabular}{@{}r r l@{}}
\toprule
$\lambda$ & $\mu_0$ for $\Li_2(\lambda)$ & Certified entries used\\
\midrule
$2/13$   & $2\,974.00$  & $31$\\
$-2/11$  & $111.27$      & $308$\\
$-2/13$  & $65.33$       & $308,328$\\
$-2/15$  & $49.22$       & $328$\\
$-3/50$  & $73\,784.43$ & $349,370$\\
$4/167$  & $70\,747.01$ & $137,147$\\
$-4/177$ & $13\,998.32$ & $137,147$\\
\bottomrule
\end{tabular}
\caption{Examples with larger reduced numerators. The last column gives zero-based entry indices in the Lean node
table. A single entry is evaluated at its node. Two entries are interpolated.}\label{tab:larger-numerators}
\end{table}

\begin{samepage}
For reduced numerator $2$ and $0<|\lambda|\le1/2$, the ranges in Corollary~\ref{cor:table} leave $2/b$ with
$b\in\{5,7,9,11\}$ and $-2/b$ with $b\in\{5,7,9\}$ outside the stated coverage. Tan and Zhang
\cite[Table~1 and \S5]{TanZhang26} give irrationality for $2/b$ from $b=13$ and for $-2/b$ from $b=17$ when
restricted to odd denominators. Liu's ranges \cite{LiuLi2Rational} give the same odd thresholds for numerator $2$.
To our knowledge the measures at $-2/11$, $-2/13$ and $-2/15$ give the first irrationality proofs for these values.
The examples $-3/50$, $4/167$ and $-4/177$ also lie below the corresponding uniform sufficient thresholds in
\cite[Table~1]{TanZhang26}.

\end{samepage}

\begin{remark}[large denominators]\label{rem:large}
For fixed $a$ the margin $\eta$ tends to infinity with $b$. Let $(\rho_k,D_k)$ be an entry of the lattice node
$\frac1{403}$. For $b\ge403|a|$ the transport with slope $m=1$ from this entry is a candidate with
$D=D_k+t-\log403$, and $c(a,b)=\frac5{12}\log b+O_a(1)$. So $\eta\ge\frac{17}{12}\log b+O_a(1)$, and
Theorem~\ref{thm:main} gives a finite measure for every large $b$ coprime to $a$. For $a=\pm1$ a numerical evaluation
of the margin for $2\le b\le403$ finds it positive, with minimum $0.1326\ldots$ at $b=2$. This remark is not part of
the formal proof.
\end{remark}

\subsection{Background}\label{ssec:background}
\emph{Classical results.} The values of polylogarithms at algebraic points are values of $G$-functions. Their
arithmetic was studied by Siegel's method and by Pad\'e approximation, and we refer to the introduction of
\cite{RhinViola05} for this history. Ap\'ery proved that $\zeta(2)=\Li_2(1)$ is irrational \cite{Apery79}, and Beukers
gave his integral proof \cite{Beukers79}. Hata constructed rational approximations to $\Li_2(1/s)$ and proved its
irrationality for every integer $s\ge7$ and $s\le-5$ \cite{Hata93dilog}. He also gave irrationality measures for
$7\le s\le18$ and $-16\le s\le-5$, from $\mu(\Li_2(1/7))<95.07$ to $\mu(\Li_2(1/18))<12.76$ and from
$\mu(\Li_2(-1/5))<228.62$ to $\mu(\Li_2(-1/16))<13.20$. Rhin and Viola extended their permutation group method from $\zeta(2)$
\cite{RhinViola96} to double integrals for the dilogarithm \cite{RhinViola05}. They proved that $\Li_2(1/6)$ is
irrational with $\mu(\Li_2(1/6))<783.29036$, and they improved Hata's measures for $7\le s\le18$. Their table gives
$\mu(\Li_2(1/s))<11.94109$ at $s=20$, and their general theorem gives $\mu(\Li_2(1/s))<16.14$ for all $s\ge15$. They
also proved that $\Li_2(r/s)$ is irrational for every $r\ge1$ and every $s$ larger than an explicit $s_1(r)$, for
instance for $r=2$ and $s\ge51$. Later they proved the linear independence of $1$, $\Li_1(r/s)$ and $\Li_2(r/s)$ for
positive and negative arguments \cite{RhinViola19}. For $-20\le s\le-5$ they gave linear independence measures $\nu$
of $1$, $\Li_1(1/s)$ and $\Li_2(1/s)$. Taking the coefficient of $\Li_1$ to be zero gives
$\mu(\Li_2(1/s))\le\nu+1$, for instance $62.69$ at $s=-5$ and $11.62$ at $s=-20$. These bounds are below Hata's for
every $-16\le s\le-5$. For $\zeta(2)$ itself the best measure is now Giraudin's $5.0193784$
\cite{Giraudin26}, after Zudilin's $5.095412$ \cite{Zudilin14}.

\emph{Small arguments.} The arguments $\pm1/2$, $\pm1/3$, $\pm1/4$ and $1/5$ are out of reach of the
classical forms. Calegari, Dimitrov and Tang recorded the irrationality of $\Li_2(1/2)$ among the open cases
\cite[Remark~2.8.2 and \S15.2]{CDT2024}. Prellberg proved that $\Li_2(-1/4)$, $\Li_2(1/5)$ and $\Li_2(-1/3)$ are
irrational by refinements of the constructions of Rhin and Viola and of Viola and Zudilin \cite{Prellberg26}. Qian Tang
adapted Fauzan's Hankel determinants for $\zeta(5)$ \cite{Fauzan2026} to the dilogarithm. His Lean formalization proves
that $\Li_2(1/2)$ is irrational \cite{TangLi2}. Ruiyi Zhang gave an independent proof \cite{ZhangLi2Half}. His
approximations are determinants of moments of $1/(xy-2)$. He bounds them by a logarithmic energy on an ellipse,
certified in interval arithmetic, and proves nonvanishing by a Lucas congruence modulo $83$.

\emph{The most recent work.} Si-Qi Liu extended Tang's method in a Lean formalization \cite{LiuLi2Rational}. He
proves that $\Li_2(1/k)$ is irrational for every $k\ge2$ and $\Li_2(-1/k)$ for every $k\ge1$. He also proves
irrationality of $\Li_2(\pm p/q)$ for $p\le4$ beyond explicit denominators, and for every $p\ge5$ when
$q\ge0.1851\,p^{84/17}$. Tan and Zhang proved that $\Li_2(p/q)$ is irrational whenever $q>(4e^4/27)p^3$, by the
Rhin--Viola integrals \cite{TanZhang26}. With Liu's construction and new comparison measures they lowered the
thresholds for $p\le7$, for instance to $q\ge418$ for $p=5$. They quote the arithmetic rate of Liu's polynomials as
$\frac83\log q-\frac94\log(q-\epsilon p)-6\log p+\frac{72}{25}$ for $\Li_2(\epsilon p/q)$ \cite[(2.2)]{TanZhang26}. This
is the rate \eqref{eq:crate} with a larger constant. These results prove irrationality. Neither states an
irrationality measure. Theorem~\ref{thm:main} uses the same polynomials and the same denominators, and it gives a
measure. As far as we know, no irrationality measure for $\Li_2(1/k)$ with $2\le k\le5$, for
$\Li_2(-1/k)$ with $2\le k\le4$, or for $\Li_2(\pm2/39)$ was published before. Table~\ref{tab:values} compares our bounds with the published ones. At every $k$ where
a published bound exists, the best published bound is that of Rhin and Viola, and ours is smaller. For $\Li_2(1/k)$
with $6\le k\le20$ the improvement runs from $783.29$ to $11.73$ at $k=6$ and from $11.94$ to $9.36$ at $k=20$. For
$\Li_2(-1/k)$ with $5\le k\le20$ it runs from $62.69$ to $10.38$ at $k=5$ and from $11.62$ to $9.30$ at $k=20$.

\emph{Recent results.} Since September 2026 several new results have appeared in this area, some obtained with AI
systems. In October 2026 OpenAI proved that Catalan's constant is irrational \cite{OpenAI2026} and that $\mu(\pi)=2$
\cite{OpenAIPi2026}, both with Lean formalizations. Liu, Jiang and Zhang proved $\mu(\log r)=2$ for rational $r>0$,
$r\ne1$ \cite{LiuJiangZhang26}. Their method needs a number that is a period of a one-dimensional group, and we do
not know how to apply it to $\Li_2(a/b)$. Fauzan proved that $\zeta(5)$ is irrational \cite{Fauzan2026}, and Lai
combined Fauzan's determinants with simultaneous approximations to show that $1$, $\zeta(3)$ and $\zeta(4)$ are
linearly independent \cite{Lai2026}.

\emph{The method.} The measure method we use, a family of integer polynomials whose values near the constant are
small and nonzero, is a one-variable form of the criteria of Nesterenko \cite{Nesterenko85} used by Ball and Rivoal
\cite{BallRivoal01}. Determinants of growing size go back to Zudilin's determinantal approach \cite{Zudilin}. The
arithmetic of the coefficients is handled prime by prime, as in the work of Hata \cite{Hata93dilog} and of Rhin and
Viola \cite{RhinViola96,RhinViola05}. The size of the determinants is controlled by Andr\'eief's identity
\cite{Forrester2018Andreief} and a logarithmic energy problem with an external field \cite{SaffTotik}.

\subsection{The companion papers}
This paper applies one method to a fifth family of constants. The method consists of Hankel or L\"owner
determinants, denominators bounded prime by prime, decay proved through an energy problem, and a formal proof of every
step. The first companion paper \cite{Litsey2026} refines OpenAI's determinant for Catalan's constant. The second
\cite{Litsey2026b} uses L\"owner determinants of digamma combinations. The third \cite{Litsey2026c} proves a measure
for $\zeta(5)$ with Fauzan's determinants, and the fourth \cite{Litsey2026d} proves measures on the plane spanned by
$\zeta(3)$ and $\pi^4$ with Lai's determinants. The dual certificate of Section~\ref{sec:decay} extends the one of
\cite{Litsey2026d}, and the comparison measures are of the arcsine type used in \cite{Litsey2026c}. We use the
definitions of $\mu$, the measure criterion and the potential-theoretic language of these papers. The construction is
Tang's, with his notation.

\subsection{What a measure requires}\label{ssec:require}
Tang's polynomials $Q_n$ have degree at most $2n$, and their values at $\xi=\Li_2(\lambda)$ are Hankel determinants of
a positive functional. Let $p/q$ be a rational close to $\xi$. Three estimates are needed at sizes $n$ of order
$\log q$.
\begin{enumerate}[label=(\roman*)]
\item \emph{Denominators.} A rational factor $d_n$ with $d_n\tilde Q_n\in\Z[X]$ primitive and $\log d_n\le-(c+o(1))n^2$,
 where $\tilde Q_n$ is Tang's normalization of $Q_n$.
 The factor removes a large common divisor of the coefficients.
\item \emph{Decay on a neighbourhood.} $|\tilde Q_n(x)|\le e^{-(D-o(1))n^2}$ for every real $x$ with
 $|x-\xi|\le e^{-(\rho+o(1))n}$, not only at $x=\xi$.
\item \emph{Nonvanishing.} $Q_n(p/q)\ne0$ for some $n$ in a range of length $O(\log q)$ above the size where (i) and
 (ii) beat $q^{2n}$.
\end{enumerate}
Then $\mu(\xi)\le\rho\max(2/(D+c),\beta)$, where $\beta$ measures how sparse the sizes with nonvanishing are. Item
(i) is known. It is the rate of Liu's construction \cite{LiuLi2Rational,TanZhang26}. The irrationality proofs of
Tang and Liu need decay at the point $\xi$ only, and nonvanishing on the line $n=\ell-1$ through the primes $\ell$. That
line has one prime per size, and it would give $\beta=1$. A measure needs the decay on a neighbourhood, for each
$a/b$, and nonvanishing at sizes of positive density.

\subsection{New ideas}
We keep Tang's determinants and the known arithmetic rate. The proof rests on the following ideas.
\begin{enumerate}[leftmargin=*]
\item \emph{Nonvanishing on a band of primes.} For every prime $\ell$ with $n/13<\ell<4n/3$, outside a finite set
 and three residues of $n$, one of $Q_n$ and $Q_{n+1}$ is nonzero at every rational whose denominator is not
 divisible by $\ell^2$ (Theorem~\ref{thm:NV}). The proof reduces the Hankel matrix modulo $\ell$ in a basis adapted to the
 residue classes of the poles. Its reduction is a product of finitely many Hankel determinants of local weights, the
 type constants, times a binomial tail in the tied case. The type constants are rational functions of $\lambda$. We
 prove that each is nonzero at every rational $\lambda$ with $0<|\lambda|\le\frac12$ by its leading terms at $0$, $1$
 and $\infty$ and one reduction modulo a prime (Theorem~\ref{thm:types}).
\item \emph{A measure criterion for a band of primes.} The primes of the band have total logarithmic mass
 $2(\frac43-\frac1{13})n=Mn$ with multiplicity two. An averaging argument over $n$ removes the exceptional residues, and
 the measure becomes $\rho\max(2/\eta,1/M)$ (Theorem~\ref{thm:criterion}). When $\eta\le2M$, nonvanishing costs nothing.
\item \emph{Decay on a neighbourhood by a capped dual.} We write the Hankel determinant at $\xi+\varepsilon$ by Heine's
 formula as a sum over configurations of $2n$ charges. The charges sit on the positive integers, at most two per site,
 or on a contour, and the perturbation adds charges at the poles $-j$, at most one per site. A positive measure
 $\sigma$ of mass $2$ and a level $\theta$ give an explicit upper bound for the rate, in which the caps enter as the
 penalties $2\int(V-2U^\sigma-\theta)_+$ and $\int(A_\rho-\theta)_+$ (Theorem~\ref{thm:LD}). For $\lambda>0$ we use the
 lattice alone. For $\lambda<0$ a bent contour through $0$ captures a cancellation that the lattice does not see.
\item \emph{Uniformity in the argument.} The dual value depends on $\lambda$ only through $t=\log(1/|\lambda|)$ and the
 sign, and it is jointly concave in the datum, $t$ and $\rho$. Lowering the level by $m\Delta$ transports a datum from
 $t$ to $t+\Delta$ at the cost $(4+m)\Delta$ in $\rho$, with gain $(2m-m^2)\Delta$. So $37$ certified data give a decay
 bound for every $\lambda$ with $|\lambda|\le\frac12$, of either sign (Theorem~\ref{thm:U}).
\item \emph{Rigorous certification of the decay bounds.} Each of the $419$ lower bounds for the dual value is proved by
 finitely many inequalities in rational arithmetic. The potentials and energies of the comparison measures have closed
 forms, and the penalty integrals are computed from the sign pattern of the field, proved on Taylor cells with explicit
 remainders (Section~\ref{ssec:cert}).
\end{enumerate}
The denominator bound of Section~\ref{sec:content} is the rate of Liu's construction. We derive it from one valuation
bound valid at every prime (Theorem~\ref{thm:K}) and Tang's estimates for the other primes. Our constant is
$\frac72-\frac{523}{840}=\frac{2417}{840}$.

\subsection{Organization}
Section~\ref{sec:construction} recalls Tang's construction. Section~\ref{sec:criterion} proves the measure criterion.
Section~\ref{sec:content} bounds the denominators, Section~\ref{sec:nv} proves the nonvanishing, Section~\ref{sec:decay}
proves the decay for one datum, and Section~\ref{sec:uniform} makes it uniform in $t$. Section~\ref{sec:proof} proves
Theorem~\ref{thm:main} and Corollary~\ref{cor:table}. Section~\ref{sec:lean} describes the formalization. The disclosure of AI
use follows it.

\subsection{Notation}
We write $v_p$ for the $p$-adic valuation and $\vG$ for the Gauss valuation of a polynomial, the minimum of the
valuations of its coefficients. We write $y_+=\max(y,0)$. Throughout, $\lambda=a/b$ is as in Theorem~\ref{thm:main},
$\xi=\Li_2(\lambda)$, $t=\log(1/|\lambda|)$ and $s=\operatorname{sign}\lambda$. For a positive measure $\sigma$ on $\C$
with compact support, $U^\sigma(z)=\int\log(1/|z-y|)\,d\sigma(y)$ is its logarithmic potential and
$I(\sigma)=\int U^\sigma d\sigma$ its energy.

%% =====================================================================
\section{Tang's construction}\label{sec:construction}

We recall the construction of \cite{TangLi2}, for a rational $\lambda$ with $0<|\lambda|<1$. Put
$D_m(x)=\prod_{i=1}^m(x+i)$ and
\[
 \tau_j=\sum_{c=1}^j\frac{\lambda^c}{c^2},\qquad \Li_{-k}(\lambda)=\sum_{x\ge1}x^k\lambda^x\in\Q .
\]
The second sum is a rational function of $\lambda$. For a variable $X$ define a $\Q$-linear functional $\Lambda_X$ on
the rational functions whose poles are simple and lie at negative integers by
\begin{equation}\label{eq:functional}
 \Lambda_X(x^k)=(k+1)\Li_{-k}(\lambda)\ \ (k\ge0),\qquad \Lambda_X\Bigl(\frac1{x+j}\Bigr)=j\lambda^{-j}(X-\tau_j)\ \ (j\ge1),
\end{equation}
extended by polynomial division and partial fractions. Its values are affine in $X$. At $X=\xi$ it has the closed
form
\begin{equation}\label{eq:sumform}
 \Lambda_\xi(R)=\sum_{x\ge1}\lambda^x\,(xR(x))' ,
\end{equation}
since $(x\cdot x^k)'=(k+1)x^k$ and $\sum_{x\ge1}\lambda^x j(x+j)^{-2}=j\lambda^{-j}(\xi-\tau_j)$. Let
\[
 W_n(x)=\frac{D_n(x)^3}{D_{4n}(x)},\qquad Q_n(X)=\det\bigl[\Lambda_X(W_n(x)\,x^{i+j})\bigr]_{0\le i,j<2n}.
\]
The entries are affine in $X$, so $\deg Q_n\le2n$. The function $W_n$ has simple poles at $-j$ for $n<j\le4n$, with
residues
\begin{equation}\label{eq:residue}
 r_j=\operatorname{Res}_{x=-j}W_n,\qquad |r_j|=\frac{((j-1)!)^2}{((j-n-1)!)^3(4n-j)!}.
\end{equation}
With $S_n=(4n)!/(n!)^3$ and $F_n=\prod_{i<2n}(i!)^2$ Tang normalizes
\[
 \tilde Q_n=\frac{S_n^{2n}}{F_n}Q_n\in\Q[X],
\]
and he writes $\tilde Q_n=\cont(\tilde Q_n)\,P_n$ with $P_n\in\Z[X]$ primitive and $\cont(\tilde Q_n)$ a positive
rational, its content. We put $\tilde d_n=1/\cont(\tilde Q_n)$. Then $P_n=\tilde d_n\tilde Q_n$, and for every rational
$p/q$ the number $q^{2n}P_n(p/q)$ is an integer.

\begin{lemma}[Heine's formula]\label{lem:heine}
For a linear functional $\Phi$ on polynomials,
\[
 \det[\Phi(x^{i+j})]_{0\le i,j<N}=\frac1{N!}\,\Phi^{\otimes N}\Bigl(\prod_{1\le k<l\le N}(x_k-x_l)^2\Bigr).
\]
\end{lemma}

\begin{proof}
Expand the square of the Vandermonde determinant $\det[x_k^i]$ and apply $\Phi$ in each variable.
\end{proof}

By \eqref{eq:sumform} the value $Q_n(\xi)$ is the Hankel determinant of the functional
$T\mapsto\sum_{x\ge1}\lambda^x(xW_nT)'(x)$. Andr\'eief's identity \cite{Forrester2018Andreief} is the same
statement for a functional given by a measure.

%% =====================================================================
\section{The measure criterion}\label{sec:criterion}

The criterion turns decay on a neighbourhood and nonvanishing on a band of primes into a measure. It is the form of
the criterion of the companion papers \cite{Litsey2026d} in which the sizes with nonvanishing are selected by primes of
positive density.

\begin{definition}\label{def:band}
Let $(P_n)_{n\ge1}$ be polynomials with rational coefficients and $0<\alpha<\beta$. A \emph{band nonvanishing datum} for
$(P_n)$ on $(\alpha,\beta)$ consists of a finite set $B$ of primes, sets $R_\ell$ of at most three residues modulo
$\ell$, and an integer $n_0$, with the following property. Let $\ell\notin B$ be a prime and $n\ge n_0$ with
$\alpha(n+1)<\ell<\beta n$ and $n\bmod\ell\notin R_\ell$. Then for all integers $p$ and $q>0$ with $\ell^2\nmid q$,
\[
 P_n(p/q)\ne0\quad\text{or}\quad P_{n+1}(p/q)\ne0 .
\]
\end{definition}

\begin{theorem}[measure criterion]\label{thm:criterion}
Let $\xi$ be real, $\kappa>0$, $\rho>0$, $\eta>0$ and $0<\alpha<\beta$. Let $P_n\in\Z[X]$ with $\deg P_n\le\kappa n$.
Suppose that for every $\varepsilon>0$ and all large $n$
\[
 |P_n(\xi+u)|\le e^{-(\eta-\varepsilon)n^2}\qquad\text{whenever } |u|\le e^{-(\rho+\varepsilon)n},
\]
and that $(P_n)$ has a band nonvanishing datum on $(\alpha,\beta)$. Put $M=2(\beta-\alpha)$. Then for every
$\mu>\rho\max(\kappa/\eta,1/M)$ there is $q_0$ with $|\xi-p/q|>q^{-\mu}$ for all integers $p$ and all $q\ge q_0$.
\end{theorem}

The proof uses three lemmas. The first is the contrapositive of Definition~\ref{def:band}.

\begin{lemma}\label{lem:pairdiv}
Let $S$ be a finite set of primes $\ell\notin B$ with $\alpha(n+1)<\ell<\beta n$, $n\ge n_0$ and $n\bmod\ell\notin R_\ell$.
If $P_n(p/q)=P_{n+1}(p/q)=0$ with $q>0$, then $\prod_{\ell\in S}\ell^2$ divides $q$.
\end{lemma}

\begin{lemma}[averaging]\label{lem:average}
For integers $N,H\ge1$ there is $n\in[N,N+H)$ with
\[
 \sum_{\substack{\ell\le2N+2H\\ n\bmod\ell\in R_\ell}}\log\ell\le3\log(2N+2H)\bigl(1+\log(2N+2H)\bigr)+\frac{6(N+H)\log4}H .
\]
\end{lemma}

\begin{proof}
For a prime $\ell$ the number of $n\in[N,N+H)$ with $n\bmod\ell\in R_\ell$ is at most $3(H/\ell+1)$. Summing
$\log\ell$ over these pairs gives at most $3H\sum_{\ell\le x}(\log\ell)/\ell+3\sum_{\ell\le x}\log\ell$ with
$x=2N+2H$. The first sum is at most $\log x\,(1+\log x)$ by the harmonic bound, and the second is at most $x\log4$ by
Chebyshev's bound. Some $n$ is at or below the average.
\end{proof}

\begin{lemma}[band mass]\label{lem:mass}
For every $\varepsilon>0$ and all large $N$ there are $n\in[N,(1+\varepsilon)N]$ and a set $S$ as in
Lemma~\ref{lem:pairdiv} with $\sum_{\ell\in S}2\log\ell\ge(M-\varepsilon)n$.
\end{lemma}

\begin{proof}
By the prime number theorem, $\sum_{\alpha(n+1)<\ell<\beta n}\log\ell\ge(\beta-\alpha-\varepsilon/4)n$ for large $n$
\cite{Davenport}. Removing the finite set $B$ costs $O(1)$. Lemma~\ref{lem:average} with $H=\lceil\varepsilon N\rceil$
gives an $n$ for which the primes with $n\bmod\ell\in R_\ell$ have mass $O(\log^2N)+O(1/\varepsilon)=o(N)$.
\end{proof}

\begin{proof}[Proof of Theorem~\ref{thm:criterion}]
Put $A=\max(\kappa/\eta,1/M)$ and fix $\mu>\rho A$. Choose $\delta,\varepsilon>0$ with $(M-\varepsilon)(1+\delta)A>1$,
$(\eta-2\varepsilon)(1+\delta)>\eta$ and $(\rho+\varepsilon)(1+\varepsilon)(1+\delta)A+\varepsilon<\mu$. Let $q$ be large,
$L=\log q$ and $N=\lceil(1+\delta)AL\rceil$. Take $n$ and $S$ from Lemma~\ref{lem:mass}, and suppose
$|\xi-p/q|\le q^{-\mu}$.

If $P_n(p/q)=P_{n+1}(p/q)=0$, then Lemma~\ref{lem:pairdiv} gives $L\ge\sum_S2\log\ell\ge(M-\varepsilon)n\ge
(M-\varepsilon)(1+\delta)AL>L$. So $P_{n'}(p/q)\ne0$ for some $n'\in\{n,n+1\}$. Then $q^{\deg P_{n'}}P_{n'}(p/q)$ is a
nonzero integer, and $|P_{n'}(p/q)|\ge q^{-\kappa n'}$. On the other hand $\mu L\ge(\rho+\varepsilon)n'$ by the choice
of $N$, so $|p/q-\xi|\le e^{-(\rho+\varepsilon)n'}$, and $|P_{n'}(p/q)|\le e^{-(\eta-\varepsilon)n'^2}$. Hence
$(\eta-2\varepsilon)n'\le\kappa L$ for large $q$. But $n'\ge(1+\delta)(\kappa/\eta)L$, so
$(\eta-2\varepsilon)(1+\delta)\le\eta$. This contradicts the choice of $\varepsilon$.
\end{proof}

The criterion concludes for all integers $p$, without coprimality, since Lemma~\ref{lem:pairdiv} and the integrality
of $q^{\deg P}P(p/q)$ hold for every fraction.

%% =====================================================================
\section{Denominators}\label{sec:content}

In this section we bound the content of $\tilde Q_n$. We have $v_p(\cont\tilde Q_n)=\vG(\tilde Q_n)$ and
\[
 \log\tilde d_n=-\sum_pv_p(\cont\tilde Q_n)\log p .
\]
So a lower bound for the Gauss valuation at each prime is what is needed.

\subsection{The local bound}
Fix a prime $p$ and put $\alpha_p=v_p(a)$, $e_p=v_p(b)$ and $f_p=v_p(b-a)$. Since $a$ and $b$ are coprime, at most one
of them is positive. Put $L_p(n)=\lfloor\log_p(7n-2)\rfloor$ and
\[
 A_2(n)=\sum_{k<2n-1}\Bigl(n+1+\Bigl\lfloor\frac k3\Bigr\rfloor\Bigr),\qquad S(n)=\sum_{i<2n}(2i-n+1)_+ .
\]

\begin{theorem}[local bound]\label{thm:K}
For every prime $p$ and every $n\ge1$,
\[
 \vG(\tilde Q_n)\ge e_pA_2(n)-\alpha_p(6n^2+n)-f_pS(n)-4nL_p(n).
\]
\end{theorem}

At a prime that divides none of $a$, $b$ and $b-a$ this is Tang's bound
$\vG(\tilde Q_n)\ge-4nL_p(n)$ \cite{TangLi2}. We describe the three other cases. Tang writes the Hankel matrix as a Gram
matrix in the binomial basis. Its entries are
\begin{equation}\label{eq:entry}
 E_{ij}=U_{ij}+\sum_{j'=1}^{4n}\rho_{j'}\binom{-j'}i\binom{-j'}j\,j'\lambda^{-j'}(X-\tau_{j'})\qquad(0\le i,j<2n),
\end{equation}
where $U_{ij}$ is the value of $\Lambda$ on the polynomial part of $S_nD_n^3\binom xi\binom xj/D_{4n}$, and $\rho_{j'}$
is the residue of $S_nD_n^3/D_{4n}$ at $-j'$. The residues $\rho_{j'}$ are $p$-integral, and they vanish for
$j'\le n$. The polynomial part enters only through the moments $\Li_{-k}(\lambda)$. In the binomial basis these are
$G(k)=r^k/(1-\lambda)$ for $k\ge1$ and $G(0)=r$, with $r=\lambda/(1-\lambda)=a/(b-a)$. So the three kinds of prime affect
three different parts of \eqref{eq:entry}.

\begin{lemma}\label{lem:V}
Let $f\ge0$ and $v_p(r)\ge-f$.
\begin{enumerate}[label=(\alph*)]
\item $v_p(G(k))\ge-(k+1)f$ for every $k\ge0$.
\item If $g$ is a polynomial of degree at most $d$ with $v_p(g(i))\ge\gamma$ for $0\le i\le d$, then the value of
 $\Lambda$ on $g$ has valuation at least $\gamma-\lfloor\log_pd\rfloor-(d+1)f$.
\item For $0\le i,j<2n$, $v_p(U_{ij})\ge-2L_p(n)-f\,(i+j-n+1)_+$.
\end{enumerate}
\end{lemma}

\begin{proof}
(a) The recurrence $G(k+1)=r(G(k)+\binom1{k+1})$ and induction. (b) Expand $g$ in the Newton basis of binomial
polynomials. The Newton coefficients have valuation at least $\gamma$, the moments lose $(d+1)f$ by (a), and the
derivative in \eqref{eq:functional} loses $\lfloor\log_pd\rfloor$. (c) If $i+j<n$, the numerator has degree below $4n$
and the polynomial part is $0$. Otherwise the polynomial part has degree at most $i+j-n$ and integral values up to
$\lfloor\log_p(3n+i+j)\rfloor\le L_p(n)$, and (b) applies.
\end{proof}

For $f=0$ these are Tang's lemmas, and the parameter $f$ is the only change.

\begin{proof}[Proof of Theorem~\ref{thm:K}]
\emph{The case $p\mid b-a$.} Then $\lambda$ and $\lambda^{-1}$ are $p$-units, and the residue part of $E_{ij}$ has Gauss
valuation at least $-2L_p(n)$ by Tang's bound for $\tau_{j'}$. With Lemma~\ref{lem:V}(c),
$\vG(E_{ij})\ge-2L_p(n)-f_p\varphi(i+j)$ with $\varphi(m)=(m-n+1)_+$. Since $\varphi$ is convex,
$\varphi(i+j)\le\frac12(\varphi(2i)+\varphi(2j))$. So $\vG(E_{ij})\ge w_i+w_j$ with $w_i=-L_p(n)-\frac{f_p}2\varphi(2i)$,
and the determinant has Gauss valuation at least $2\sum_iw_i=-4nL_p(n)-f_pS(n)$.

\emph{The case $p\mid b$.} Then $r$ is a $p$-unit and $v_p(\lambda^{-1})=e_p\ge1$. This is Tang's argument for $p=2$ at
$\lambda=\frac12$, with $e_p$ in place of $1$ \cite{TangLi2}. The Hankel matrix factors as a product of a matrix of
values at the poles and a matrix of binomial coefficients. Tang attaches a weight to each of the $2n$ slots of the
factorization, ranks the slots by their weights, and applies the Cauchy--Binet expansion. The weights of the pole slots
are $e_p\min(n+1+\lfloor(i-1)/3\rfloor,4n+1)-2L_p(n)$, and their sum over the $2n$ ranks is $e_pA_2(n)-4nL_p(n)$.

\emph{The case $p\mid a$.} Then $r$ and $\lambda$ have nonnegative valuation, and $v_p(\lambda^{-j'})=-\alpha_pj'$. By
\eqref{eq:entry} the matrix $[E_{ij}]$ is the product $WC$ of a $2n\times6n$ matrix $W$ and a $6n\times2n$ matrix $C$.
The first $2n$ columns of $W$ are the polynomial parts $U_{i\cdot}$, and the others are
$\rho_{j'}\binom{-j'}ij'\lambda^{-j'}(X-\tau_{j'})$ for $1\le j'\le4n$. The rows of $C$ are the unit vectors and the
integers $\binom{-j'}j$. The slots have weights $-2L_p(n)$ and $-2L_p(n)-\alpha_pj'$. Rank the slot $j'$ as $4n-j'$ and
the polynomial slots above them. By the Cauchy--Binet expansion each minor of $W$ uses $2n$ distinct slots, so its
weight is at least $\sum_{t<2n}(-2L_p(n)-\alpha_p(4n-t))=-4nL_p(n)-\alpha_p(6n^2+n)$.
\end{proof}

\subsection{The rate}
The other primes are handled by Tang's estimates \cite{TangLi2}, which hold for every rational $\lambda$ with
$|\lambda|<1$ at the primes that divide none of $a$, $b$ and $b-a$. We state them as we use them.

\begin{proposition}[Tang]\label{prop:tang}
Let $\varepsilon>0$ and $\delta=1/(N+1)$ with $4/(N+1)<\varepsilon/2$. For large $n$, with the sums over primes $p$ not
dividing $ab(b-a)$:
\begin{enumerate}[label=(\alph*)]
\item (small primes) $\sum_{4<p\le\delta n}v_p(\cont\tilde Q_n)\log p\ge(-4\delta-\varepsilon/8)n^2$;
\item (medium primes) $\sum_{\delta n<p\le n}v_p(\cont\tilde Q_n)\log p\ge(-\frac{523}{840}+\frac23\delta^2-\varepsilon/8)n^2$;
\item (outer primes) $\sum_{n<p\le4n}v_p(\cont\tilde Q_n)\log p\ge(\frac72-\varepsilon/8)n^2$;
\item (large primes) $v_p(\cont\tilde Q_n)\ge0$ for $p>4n$.
\end{enumerate}
\end{proposition}

Part (a) is the local bound summed with Chebyshev's estimate. In part (b) Tang bounds $v_p$ by a profile
$n\,\phi(p/n)$, where $\phi(x)$ is an explicit function of $\lfloor1/x\rfloor$, $\lfloor2/x\rfloor$ and $\lfloor4/x\rfloor$.
The weighted prime number theorem turns the sum into $n^2\int_\delta^1\phi$. On each interval where the floors are
constant the integral is a rational function of the floor values, and the sum of these is bounded below by
$-523/840$. Part (c) counts the multiples of $p$ among the factors of $S_n^{2n}/F_n$ and the poles. Part (d) holds
because every factor is then $p$-integral.

\begin{theorem}[arithmetic rate]\label{thm:content}
For every $\varepsilon>0$ there is $n_1$ such that $\log\tilde d_n\le-(c(a,b)-\varepsilon)n^2$ for $n\ge n_1$ with
$\tilde Q_n\ne0$.
\end{theorem}

\begin{proof}
Split the primes $p\le4n$ into the set $\mathcal S$ of primes dividing $ab(b-a)$, the primes $p\le4$, and the three
ranges of Proposition~\ref{prop:tang}. For $n$ large the last two ranges miss $\mathcal S$. By
Theorem~\ref{thm:K}, $A_2(n)\ge\frac83n^2-O(n)$, $S(n)\le\frac94n^2$ and $L_p(n)=O(\log n)$. So
\[
 \sum_{p\in\mathcal S}v_p(\cont\tilde Q_n)\log p\ge\Bigl(\frac83\log b-6\log|a|-\frac94\log|b-a|-\frac\varepsilon8\Bigr)n^2 .
\]
The primes $p\le4$ outside $\mathcal S$ cost $O(n\log n)$. Adding Proposition~\ref{prop:tang} and letting $\delta\to0$
with $\varepsilon$ gives the claim.
\end{proof}

\begin{remark}
The rate of Theorem~\ref{thm:content} is the rate of Liu's construction, quoted by Tan and Zhang
\cite[(2.2)]{TanZhang26}. The logarithmic terms are the same. Liu's constant is $\frac{72}{25}$, and ours is
$\frac{2417}{840}=\frac{72}{25}-\frac{11}{4200}$. Our constant comes from Tang's estimate for the medium primes in
Proposition~\ref{prop:tang}(b).
\end{remark}

%% =====================================================================
\section{Nonvanishing}\label{sec:nv}

\subsection{Statement}
Let $\ell\ge13$ be a prime. Write $n=q\ell+r$ with $0\le r<\ell$, $k=\lfloor4r/\ell\rfloor$ and $r'=4r-k\ell$. We call
$n$ \emph{tied} at $\ell$ if $r'<r$. We assume
\begin{equation}\label{eq:range}
 \ell(\ell-1)>4n,\qquad\text{and}\qquad \ell\le n\ \text{or}\ 3\ell<4n<4\ell .
\end{equation}
For integers $0\le\alpha'<\gamma'$ and $m\ge0$ define the \emph{local weights}
\[
 w_{\alpha',\gamma'}(u)=\frac{\prod_{m'<\alpha'}(u+m')^2}{\prod_{m'=\alpha'}^{\gamma'-1}(u+m')},\qquad
 w^0_{\alpha',\gamma'}(u)=\frac{\prod_{m'=1}^{\alpha'}(u+m')^2}{\prod_{m'=\alpha'+1}^{\gamma'}(u+m')},
\]
and the \emph{type constants}
\begin{gather*}
 K_V(\alpha',\gamma';m)=\det\bigl[\mathcal V(u^{i+i'}w_{\alpha',\gamma'})\bigr]_{i,i'<m},\\
 K_U(\alpha',\gamma';m)=\det\bigl[\mathcal V(u^{i+i'+1}w^0_{\alpha',\gamma'})\bigr]_{i,i'<m},
\end{gather*}
where $\mathcal V$ is the functional $f\mapsto\sum_{x\ge1}\lambda^xf'(x)$ with the coefficient of $\Li_2(\lambda)$ set to
$0$. Concretely $\mathcal V(u^j)=j\Li_{1-j}(\lambda)$ and $\mathcal V(1/(u+m))=\lambda^{-m}\tau_m$. These are
rational functions of $\lambda$. For $q\le12$ and $k\le3$ let $\mathcal K(q,k)$ be the finite list of constants
attached to the residue types below, together with $K_V(q+1,4q+k;2q+k-1)$ and $K_U(q,4q+k;2q+k-1)$ when $k\ge1$.

\begin{theorem}[nonvanishing]\label{thm:NV}
Let $\ell$ and $n$ satisfy \eqref{eq:range} with $q\le12$, and let $\ell$ divide none of the numerators and
denominators of $\lambda$ and $1-\lambda$. Suppose that every constant of $\mathcal K(q,k)$ is an $\ell$-adic unit and,
if $n$ is tied, that the number $T'(n,\ell)$ defined in \eqref{eq:Tprime} is not divisible by $\ell$. Then
$Q_n(x)\ne0$ for every $x\in\Q$ with $v_\ell(x)\ge-1$.
\end{theorem}

The tolerance $v_\ell(x)\ge-1$ allows a rational $p/q$ whose denominator is divisible by $\ell$ but not by $\ell^2$.

\subsection{Residue classes}
For $0\le\sigma<\ell$ let $a_\sigma$ be the number of $i\in[1,n]$ and $c_\sigma$ the number of $j\in[1,4n]$ with
$i\equiv\sigma$, $j\equiv\sigma\pmod\ell$. Put $m_\sigma=c_\sigma-2a_\sigma$, $b_\sigma=c_\sigma-a_\sigma$ and
$\nu_\sigma=3a_\sigma-c_\sigma-1+[\sigma=0]$. Then $\sum_\sigma m_\sigma=2n$. Under \eqref{eq:range} we have
$c_\sigma\le\ell-1$ and $m_\sigma\ge0$, and the classes fall into five types.

\begin{center}
\begin{tabular}{@{}lccc@{}}
\toprule
type & classes & $(a_\sigma,c_\sigma)$ & $m_\sigma$\\
\midrule
A & $1\le\sigma\le\min(r,r')$ & $(q+1,4q+k+1)$ & $2q+k-1$\\
B & $r<\sigma\le r'$ & $(q,4q+k+1)$ & $2q+k+1$\\
C & $r'<\sigma\le r$ & $(q+1,4q+k)$ & $2q+k-2$\\
D & $\sigma>\max(r,r')$ & $(q,4q+k)$ & $2q+k$\\
zero & $\sigma=0$ & $(q,4q+k)$ & $2q+k$\\
\bottomrule
\end{tabular}
\end{center}

The types B and C never occur together, and C occurs exactly when $n$ is tied. The type constant of a class of type A,
B, C or D is $K_V(a_\sigma,c_\sigma;m_\sigma)$, and that of the zero class is $K_U(a_0,c_0;m_0)$.

\subsection{Proof of Theorem~\ref{thm:NV}}
We follow the reduction of the Hankel matrix modulo $\ell$ in a basis adapted to the classes.

\emph{Step 1: the basis.} Put $\Pi_\sigma=\prod_{\sigma'\ne\sigma}(x+\sigma')^{m_{\sigma'}}$ and
$e_{\sigma,i}=(x+\sigma)^i\Pi_\sigma$ for $i<m_\sigma$. These are $2n$ polynomials of degree below $2n$. Their
coefficient matrix $T$ has determinant prime to $\ell$, by the Chinese remainder theorem, since the roots
$-\sigma$ are distinct modulo $\ell$. The Gram matrix $G(x)=[\Lambda_x(W_ne_\alpha e_\beta)]$ satisfies
$\det G(x)=(\det T)^2Q_n(x)$. So it suffices to show $\det G(x)\ne0$.

\emph{Step 2: the $\ell$-adic functional.} Put $\lambda_\ell=\lambda^\ell$, so that $\lambda_\ell\equiv\lambda\pmod\ell$,
and note that $\Li_{-k}(\lambda_\ell)$ is $\ell$-integral since $1-\lambda_\ell$ is a unit. For a power series
$h=\sum H_jy^j$ in $y=\ell u$ with $\ell$-integral coefficients put $\mathcal S(h)=\sum_jH_j\ell^j\Li_{-j}(\lambda_\ell)$,
a convergent $\ell$-adic sum. The identity $\lambda\sum_{i\le j}\binom ji\Li_{-i}(\lambda)=\Li_{-j}(\lambda)-\lambda$ gives
the shift rule $\mathcal S(h(\cdot+1))=\lambda_\ell^{-1}(\mathcal S(h)-\lambda_\ell h(1))$. For $\ell\nmid c$ put
$E(c)=\mathcal S((c+\ell u)^{-2})$. The shift rule gives $E(c+\ell)=\lambda_\ell^{-1}E(c)-(c+\ell)^{-2}$, and from this
the sum
\[
 T(j)=\sum_{\substack{c\le j\\ \ell\nmid c}}\frac{\lambda^c}{c^2}+\sum_{\substack{0\le\sigma<\ell\\ \sigma\not\equiv j}}\lambda^{j-\sigma}E(j-\sigma)
\]
does not depend on $j$. Call its value $L$. It is $\ell$-integral.

\emph{Step 3: decomposition by classes.} For a rational function $F$ with poles at negative integers and a class
$\sigma$, expand $(xF)'$ at $x=-\sigma+\ell u$, and let $\mathcal W_\sigma(F;x)$ be $\lambda^{-\sigma}$ times the
corresponding $\ell$-adic sum, with $\ell^2(x-L)$ in place of the value $\Li_2(\lambda_\ell)$. Then
$\Lambda_x(F)=\sum_{\sigma<\ell}\mathcal W_\sigma(F;x)$. For $F=x^k$ this is the distribution relation
$\Li_{-k}(\lambda)=\sum_\sigma\lambda^{-\sigma}\sum_i\binom ki(-\sigma)^{k-i}\ell^i\Li_{-i}(\lambda_\ell)$, an identity of
rational functions. For $F=1/(x+j)$ it follows from the shift rule and the constancy of $T$. Since $v_\ell(L)\ge0$,
$x-L$ has valuation at least $-1$ when $x$ does.

\emph{Step 4: local expansion.} At $x=-\sigma+\ell u$ a factor $x+i$ with $i\equiv\sigma$ equals $\ell(u+m')$, and the
other factors are $\ell$-adic units times $1+O(\ell u)$. So $W_n$ is $\ell^{3a_\sigma-c_\sigma}$ times a unit times
$w_\sigma(u)$ times a power series. It follows that $\mathcal W_\sigma(W_ne_\alpha e_\beta;x)$ has valuation at least
$\nu_\sigma+d_\sigma(\alpha)+d_\sigma(\beta)$, where $d_\sigma(\alpha)$ is the order of vanishing of $e_\alpha$ at $-\sigma$
modulo $\ell$. The part proportional to $\ell^2(x-L)$ has one more power of $\ell$. This is where the tolerance enters.
For two basis vectors of the same class $\sigma$ the leading term is a unit times
$\mathcal V(u^{i+i'+[\sigma=0]}w_\sigma)$ reduced modulo $\ell$.

\emph{Step 5: thresholds.} Put $\pi_{\sigma,i}=i+\nu_\sigma/2$ and $K_\sigma=m_\sigma+\nu_\sigma/2$. By the table,
all $K_\sigma$ lie in an interval of length one. A comparison of
$\nu_{\sigma''}+d_{\sigma''}(\alpha)+d_{\sigma''}(\beta)$ with $\pi_\alpha+\pi_\beta$ shows that the entries of the
scaled matrix $Z_{\alpha\beta}=\ell^{-\pi_\alpha-\pi_\beta}G_{\alpha\beta}$, with half-integers split between rows and
columns by the parity of $\nu_\sigma$, are $\ell$-integral. Modulo $\ell$ the cross-class entries vanish, except in
the tied case, where the classes of type C couple to the top row of the zero class.

\emph{Step 6: the reduction.} If $n$ is not tied, $Z\bmod\ell$ is block diagonal over the classes, and its
determinant is a product of units times the type constants. If $n$ is tied, the group formed by the classes of type C
and the zero class has a bordered block structure. Its Schur complement gives
\begin{equation}\label{eq:Tprime}
 T'(n,\ell)=c_1-c_2B(r,r'),\qquad B(r,r')=\sum_{i=r'+1}^r\binom ri\lambda^{-i}(1-\lambda^{-1})^{r-i},
\end{equation}
times a unit, where $c_1=K_U(q,4q+k;2q+k)K_V(q+1,4q+k;2q+k-2)$ and $c_2=K_U(q,4q+k;2q+k-1)K_V(q+1,4q+k;2q+k-1)$. The
binomial tail comes from the units of Step~4, which Wilson's theorem evaluates modulo $\ell$, and from the polynomial
identity
\[
 \sum_{i=r'+1}^r\binom riy^i(1-y)^{r-i}=\frac{r!}{r'!\,N!}\sum_{j=0}^N(-1)^j\binom Nj\frac{y^{r'+1+j}}{r'+1+j}\qquad(N=r-r'-1).
\]
In both cases $\det Z\not\equiv0\pmod\ell$, so $\det G(x)\ne0$.

\subsection{The pair rule}
When both $n$ and $n+1$ are tied, a second congruence comes from Pascal's rule. The residue $k$ is the same, and
$r'(n+1)=r'(n)+4$. If $T'(n,\ell)\equiv T'(n+1,\ell)\equiv0$ and $\ell\nmid c_2$, then $\ell$ divides
$R_y(r)$, where $y=\lambda^{-1}$ and
\begin{multline*}
 R_y(r)=4(1-y)^3y(4r+1)(4r+2)(4r+3)-4(1-y)^3y(r+1)(4r+2)(4r+3)\\
 +12(1-y)^2y^2r(r+1)(4r+3)-12(1-y)y^3r(r+1)(3r+1)\\
 +3y^4r(3r+1)(3r+2).
\end{multline*}
As a polynomial in $r$, $R_y=r(\rho_1(y)+\rho_2(y)r+\rho_3(y)r^2)$ with
$\rho_3(y)=-3y(y-4)(27y^2-40y+16)$, and the gcd of $\rho_1,\rho_2,\rho_3$ in $\Z[y]$ is $3y$. So for every rational
$y\ne0$ the polynomial $R_y$ is nonzero, and it has at most three roots modulo $\ell$ for all $\ell$ outside a finite
set.

\begin{corollary}\label{cor:band}
Suppose that every type constant of $\mathcal K(q,k)$, $q\le12$, $k\le3$, is a nonzero rational at $\lambda$. Then the
family $(P_n)$ has a band nonvanishing datum on $(\frac1{13},\frac43)$.
\end{corollary}

\begin{proof}
Let $B$ be the primes below $13$, the primes dividing the numerators and denominators of $6\lambda(1-\lambda)$, of the
constants of $\mathcal K(q,k)$, of $c_2$ and of a nonzero coefficient of $R_y$. Let $R_\ell$ be the roots of $R_y$
modulo $\ell$. If $\frac{n+1}{13}<\ell<\frac43n$ and $n$ is large, then $n$ and $n+1$ satisfy \eqref{eq:range} with
$q\le12$. If $n$ is not tied, or $T'(n,\ell)\not\equiv0$, Theorem~\ref{thm:NV} applies at $n$. Otherwise it applies at
$n+1$, or both are exceptional and $n\bmod\ell\in R_\ell$. Since $P_n$ is a nonzero rational multiple of $Q_n$, the
conclusion holds for $P_n$.
\end{proof}

\subsection{The type constants}
It remains to show that the constants are nonzero. There are $318$ of them for $q\le12$. Each is a rational function
of $\lambda$, and $K_V(1,4;2)=-(\lambda+12)(13\lambda-12)/(8\lambda^3(\lambda-1))$ shows that some vanish at rational
points. Let $Z_0=\{-12,-4,-\frac45,\frac{12}{13}\}$.

\begin{theorem}\label{thm:types}
Every constant of $\mathcal K(q,k)$ with $q\le12$ and $k\le3$ is nonzero at every rational $\lambda\notin
Z_0\cup\{0,1\}$. In particular it is nonzero for $0<|\lambda|\le\frac12$.
\end{theorem}

The proof rests on a principle for rational roots and on the leading terms of the constants at $0$, $1$ and $\infty$.

\begin{lemma}\label{lem:R}
Let $F\in\Z[z]$ be nonzero and $p$ prime. Let $v_0$ and $v_1$ be the orders of $F$ at $0$ and $1$. Suppose that $p$ does
not divide the lowest coefficient of $F$ at $0$, the lowest Taylor coefficient of $F$ at $1$, or the leading
coefficient of $F$, and that $F\bmod p$ has no root in $\F_p\setminus\{0,1\}$. Then $F(z)\ne0$ for every
$z\in\Q\setminus\{0,1\}$.
\end{lemma}

\begin{proof}
Let $z=a'/b'$ in lowest terms and $G=b'^{\deg F}F(a'/b')$. If $p\nmid a'b'(b'-a')$, then $G\not\equiv0$ by the root
condition. If $p\mid a'$, then $G=a'^{v_0}(f_{v_0}b'^{\deg F-v_0}+a'(\dots))$ and the bracket is prime to $p$. The cases
$p\mid b'-a'$ and $p\mid b'$ are the same with the Taylor expansion at $1$ and the leading coefficient.
\end{proof}

\begin{lemma}\label{lem:potentials}
Let $M$ be an $m\times m$ matrix over the Laurent series in one variable $x$, and let $r_i,c_j$ be integers with
$\operatorname{ord}M_{ij}\ge r_i+c_j$. Let $\Lambda_{ij}$ be the coefficient of $x^{r_i+c_j}$ in $M_{ij}$. Then
$\det M=x^{\sum r_i+\sum c_j}(\det\Lambda+O(x))$.
\end{lemma}

\begin{proof}
Multiply by $\operatorname{diag}(x^{-r_i})$ and $\operatorname{diag}(x^{-c_j})$ and reduce to the constant term.
\end{proof}

\begin{proof}[Proof of Theorem~\ref{thm:types}]
Write the weight $w$ of a constant as a polynomial part plus $\sum_m\varrho_m/(u+m)$. Then
\[
 \mathcal V(g\,w)=\sum_m\varrho_mg(-m)\lambda^{-m}\tau_m+\sum_{k\ge1}k\,p_k\Li_{1-k}(\lambda),
\]
with $p_k$ the coefficients of the polynomial part of $gw$. At each of the points $\lambda=0$, $1$, $\infty$ we choose
a basis in which the Hankel matrix has an explicit pattern of orders, and Lemma~\ref{lem:potentials} gives the leading
coefficient.
\begin{itemize}
\item At $\lambda=0$ we use the Newton basis on the largest poles $m_1>m_2>\dots$. Since $\lambda^{-m}\tau_m$ has order
 $1-m$ with leading coefficient $1$, the matrix of leading terms is triangular. The leading coefficient is
 $\prod_k\varrho_{m_k}\prod_{k<l}(m_k-m_l)^2$.
\item At $\lambda=1$ the pole part is regular with value $\sum_m\varrho_mg(-m)H^{(2)}_m$, where $H^{(2)}_m=\sum_{c\le m}c^{-2}$,
 and $\Li_{-j}(\lambda)$ has a pole of order $j+1$ with leading coefficient $j!$. The matrix of leading terms is block
 triangular. Its determinant is one Hankel determinant of size at most $8$ in the numbers
 $\sum_m\varrho_m(-m)^iH^{(2)}_m$, times products of factorials. This uses
 $\det[(i+j+\gamma)!]_{i,j<N}=\prod_{k<N}k!(k+\gamma)!$.
\item At $\lambda=\infty$ we use $\Li_{-j}(1/w)=(-1)^{j+1}\Li_{-j}(w)$ and the Hermite--Newton basis on the $3$-jets of
 the poles. The matrix of leading terms is block triangular with anti-triangular blocks, of determinant
 $\varrho_m^3$ times squares of values of the basis, and one block of size at most $2$.
\end{itemize}
For each constant let $F\in\Z[\lambda]$ be the numerator obtained by clearing the denominators of the moments. The
three leading coefficients of $F$ are the closed forms above, up to a power of a fixed integer. For each of $313$ of the
$318$ constants we fix a prime $p$ and check in exact arithmetic that $p$ divides none of the three leading coefficients and
that $F\bmod p$ has no root in $\F_p\setminus\{0,1\}$. Lemma~\ref{lem:R} then gives the claim. The other five constants have
$q\le1$. For them $F$ has degree at most $10$, and an exact factorization shows that its rational roots lie
in $Z_0\cup\{0,1\}$.
\end{proof}

%% =====================================================================
\section{Decay}\label{sec:decay}

\subsection{The field and the dual value}
Put
\[
 \Phi_0(z)=3\int_0^1\log|z+y|\,dy-\int_0^4\log|z+y|\,dy .
\]
This is the limit of $n^{-1}\log|W_n(nz)|+\log n$. For $1\le x\le4$,
$\Phi_0(-x)=2x\log x-3(x-1)\log(x-1)-(4-x)\log(4-x)+1$. Let $\kappa_s=2\pi$ for $s=+1$ and $\kappa_s=\pi$ for $s=-1$. Define
on $\{\operatorname{Re}z\ge0\}$ and on $[1,4]$
\[
 V(z)=\Phi_0(z)-t\operatorname{Re}z-\kappa_s|\operatorname{Im}z|,\qquad \omega(x)=\Phi_0(-x)+tx .
\]
The first is the rate of the weights $\lambda^x$ and of the kernel of a contour. The second is the rate of the
residues $r_jj\lambda^{-j}$ at $j=xn$, by \eqref{eq:residue} and Stirling's formula.

We use two kinds of frame. The \emph{lattice frame} has no contour, and every integer $x\ge1$ is a site. A \emph{bent
frame} is used for $\lambda<0$. It is a polyline $\Gamma$ in $\{\operatorname{Re}z\ge0\}$ from $-i\infty$ to $+i\infty$,
symmetric under conjugation, that meets $\R$ only at $0$, has no horizontal edge through $0$, and ends in vertical
rays. We use four bent frames $\Gamma_2,\Gamma_3,\Gamma_4,\Gamma_{15}$. Their upper vertices after $0$ are
\begin{gather*}
 \Gamma_2:\ \tfrac1{200}+\tfrac18i,\ \tfrac{89}{100}+\tfrac{457}{200}i;\qquad
 \Gamma_3:\ \tfrac25+\tfrac{18}{25}i,\ \tfrac75+\tfrac{68}{25}i;\\
 \Gamma_4:\ \tfrac6{25}+\tfrac3{10}i,\ \tfrac{167}{100}+\tfrac{263}{100}i;\qquad
 \Gamma_{15}:\ \tfrac{13}{50}+\tfrac3{50}i,\ 2+2i ,
\end{gather*}
each followed by a vertical ray.

A \emph{dual datum} is a pair $(\sigma,\theta)$ with $\theta\in\R$ and $\sigma$ a positive measure of mass $2$ and
compact support whose potential is H\"older continuous. Finite sums of measures on segments with bounded density
against the uniform or the arcsine law have this property. Put
\[
 W_\sigma=V-2U^\sigma,\qquad A_\rho(x)=\omega(x)-\rho-2U^\sigma(-x)\quad(1\le x\le4).
\]
The datum is \emph{admissible} if $W_\sigma\le\theta$ on $\Gamma$. For the lattice frame there is no condition. The
\emph{dual value} is
\begin{multline}\label{eq:dual}
 D_\rho(\sigma,\theta)=-(12\log2+4)-I(\sigma)-2\theta\\
 -2\int_0^\infty(W_\sigma(y)-\theta)_+\,dy-\int_1^4(A_\rho(x)-\theta)_+\,dx
\end{multline}
for the lattice frame, and the same without the lattice integral for a bent frame.

\begin{theorem}[decay]\label{thm:LD}
Let $(\sigma,\theta)$ be admissible for the lattice frame, or for a bent frame when $\lambda<0$. Then for every
$\rho>0$
\[
 \limsup_{n\to\infty}\frac1{n^2}\log\sup_{|\varepsilon|\le e^{-\rho n}}|\tilde Q_n(\xi+\varepsilon)|\le-D_\rho(\sigma,\theta).
\]
\end{theorem}

The dual value depends on $\lambda$ only through $t$ and $s$. For the lattice frame it does not depend on $s$.

\subsection{Heine's formula with jets}
By \eqref{eq:functional} and \eqref{eq:sumform},
$\Lambda_{\xi+\varepsilon}(W_nT)=\Lambda_\xi(W_nT)+\varepsilon\sum_{j=n+1}^{4n}r_jj\lambda^{-j}T(-j)$. For the lattice frame
$\Lambda_\xi(W_nT)=\sum_{x\ge1}(\alpha_xT(x)+\beta_xT'(x))$ with $\alpha_x=\lambda^x(xW_n)'(x)$ and
$\beta_x=\lambda^xxW_n(x)$. So the functional is a sum of jets at the sites and point masses at the poles.

\begin{lemma}\label{lem:cap}
In Heine's formula for this functional, a configuration contributes only if every site carries at most two variables.
A site with one variable contributes $(\alpha_x+\beta_x\partial)$ applied to the square of the Vandermonde product, and
a site with two variables $u_k,u_l$ contributes $-2\beta_x^2$ times the product with the factor $(u_k-u_l)^2$ removed.
\end{lemma}

\begin{proof}
At a site carrying $m$ variables the operator has order at most $m$. The square of the Vandermonde product vanishes
to order $m(m-1)$ on the diagonal, which exceeds $m$ for $m\ge3$. For $m=2$ only $\beta_x^2\partial_k\partial_l$ survives,
and $\partial_k\partial_l(u_k-u_l)^2=-2$.
\end{proof}

A derivative brings at most a factor $Cn$, since $|\alpha_x/\beta_x|\le5(1+\log n)$ and the other variables are at
distance at least $\frac12$. So every term is bounded by $e^{O(n\log n)}$ times the product of the weights of its
charges and the reduced Vandermonde product.

\subsection{The contour}
For a bent frame the sum over the sites is replaced by a contour integral.

\begin{lemma}\label{lem:contour}
Let $\lambda<0$, $L=\log|\lambda|$ and $T$ a polynomial. Let $\Gamma_n$ be $n\Gamma$ with its part in
$|\operatorname{Im}z|\le1$ replaced by the segment $\{\frac12+iy:|y|\le1\}$ and two horizontal joins. Then
\[
 \sum_{x\ge1}\lambda^x(xW_nT)'(x)=\frac1{2\pi i}\int_{\Gamma_n}T(z)\,zW_n(z)|\lambda|^z\,\kappa(z)\,dz,
\]
where $\kappa=-k'-Lk$ and $k(z)=\pi/\sin\pi z$,
with $\Gamma_n$ oriented downward, and $|\kappa(z)|\le Ce^{-\pi|\operatorname{Im}z|}$ on $\Gamma_n$.
\end{lemma}

\begin{proof}
The kernel $k$ has residue $(-1)^x$ at $x$. The residue theorem on the region to the right of $\Gamma_n$ gives the sum
with $(H'-LH)k$, $H(z)=zW_n(z)T(z)|\lambda|^z$, since the arcs at infinity contribute $0$. An integration by parts
moves the derivative to the kernel. On the vertical segment $|\sin\pi z|\ge1$, on the joins $|\sin\pi z|\ge\sinh\pi$,
and for $|\operatorname{Im}z|\ge1$ the bound follows from $|\sin\pi z|\ge\frac12e^{\pi|\operatorname{Im}z|}(1-e^{-2\pi})$.
\end{proof}

\subsection{The weights}
\begin{lemma}\label{lem:weights}
(a) For $\operatorname{Re}u\ge0$ and $|u|\ge\frac12$,
$\log|W_n(u)|\le n\Phi_0(u/n)-n\log n+3\log(|u|+n+1)+3\log2$.
(b) For $n<j\le4n$, $\log|r_jj\lambda^{-j}|=n\,\omega(j/n)-n\log n+O(\log n)$.
\end{lemma}

\begin{proof}
(a) The function $y\mapsto\log|u+y|$ is nondecreasing on $y\ge0$. So the sum over $1\le i\le n$ is at most the
integral over $[1,n+1]$, and the sum over $1\le j\le4n$ is at least the integral over $[0,4n]$. Substitute $y=nx$.
(b) Stirling's formula in \eqref{eq:residue}.
\end{proof}

\subsection{The dual inequality}
Scale every charge by $n$. A configuration then consists of $2n$ points $s_i$: lattice points $x/n$ with at most two
per site, contour points, and points $-j/n$ with at most one per pole.

\begin{lemma}\label{lem:dual}
Let $\sigma$ be a positive measure of mass $2$ and compact support with a H\"older continuous potential. For every
configuration of $2n$ points,
\[
 \sum_{i\ne j}\log|s_i-s_j|\le n^2I(\sigma)-2n\sum_iU^\sigma(s_i)+o(n^2),
\]
uniformly in the configuration, where a confluent pair counts with $\log2$ in place of $\log0$.
\end{lemma}

\begin{proof}
Replace each point by the normalized arc measure on a circle of radius $n^{-2}$ around it, and let $\mu^\epsilon$ be
$n^{-1}$ times their sum. Subharmonicity bounds $\log|s_i-s_j|$ by the average over the two circles. The energy of
the neutral measure $\mu^\epsilon-\sigma$ is nonnegative \cite[Lemma~I.1.8]{SaffTotik}, so
$-I(\mu^\epsilon)\le I(\sigma)-2\int U^\sigma d\mu^\epsilon$. The H\"older continuity of $U^\sigma$ replaces the averages
over the circles by the values at the points with an error $o(n^2)$ in total.
\end{proof}

\begin{proof}[Proof of Theorem~\ref{thm:LD}]
By Lemma~\ref{lem:heine} and Lemma~\ref{lem:cap}, $|Q_n(\xi+\varepsilon)|$ is at most $e^{O(n\log n)}$ times a sum over
configurations of the product of the weights and of the reduced Vandermonde product. By Lemma~\ref{lem:weights} and
Lemma~\ref{lem:contour}, a lattice or contour charge at $s$ has weight at most $e^{n(V(s)-\log n)+O(\log n)}$, and an
atom charge at $-x$ has weight at most $|\varepsilon|e^{n(\omega(x)-\log n)+O(\log n)}\le e^{n(\omega(x)-\rho-\log n)+O(\log n)}$.
Put $f(s)=W_\sigma(s)$ for a lattice or contour charge and $f(-x)=A_\rho(x)$ for an atom charge. The Vandermonde
product contributes $n^{2n(2n-1)}=e^{4n^2\log n+O(n\log n)}$ from the scaling, and the $2n$ weights contribute
$e^{-2n^2\log n}$. By Lemma~\ref{lem:dual} the logarithm of a term is at most
\[
 n^2I(\sigma)+n\sum_if(s_i)+2n^2\log n+o(n^2).
\]
For any $\theta$,
\[
 \sum_if(s_i)\le2n\theta+\sum_i(f(s_i)-\theta)_+ .
\]
Contour charges have $W_\sigma\le\theta$ by admissibility. Lattice charges occupy each site at most twice, so their
excess is at most $2n\int_0^\infty(W_\sigma-\theta)_++o(n)$. Atom charges occupy each pole at most once, so theirs is
at most $n\int_1^4(A_\rho-\theta)_++o(n)$. The number of configurations is $e^{O(n\log n)}$ after the tails are summed,
since $W_\sigma(z)\le3\log(|z|+R)-t\operatorname{Re}z-\kappa_s|\operatorname{Im}z|+O(1)$ for large $|z|$. Finally
$\log(S_n^{2n}/F_n)=-2n^2\log n+(12\log2+4)n^2+O(n\log n)$, which cancels the power of $n$ and contributes the
constant. Adding these bounds gives $\log|\tilde Q_n(\xi+\varepsilon)|\le-D_\rho(\sigma,\theta)n^2+o(n^2)$.
\end{proof}

\subsection{Certificates}\label{ssec:cert}
A dual datum gives a lower bound for the decay rate once three quantities are bounded: the energy $I(\sigma)$, the
admissibility on $\Gamma$, and the two penalty integrals. Our data are of two shapes.

For the lattice frame, $\sigma$ is the sum of a measure on $[0,Y]$ and a measure on $[-4,-1]$, the second carrying the
share of the poles. Each is the arcsine law of its segment times $|\sum_kh_ke^{ik\phi}|^2$, with $t=\cos\phi$ and
rational $h_k$. Its potential inside and outside the segment and its energy are explicit Chebyshev series, as in
\cite{Litsey2026c}. The cross energy of the two segments is a sum of Stieltjes moments. The penalty
$\int(W_\sigma-\theta)_+$ is computed from the sign pattern of $W_\sigma-\theta$. Between consecutive rational points
where the sign is constant, the integral of $(W_\sigma-\theta)_+$ is either $0$ or an explicit integral of
$W_\sigma-\theta$. The signs are proved on cells by Taylor polynomials with explicit remainders. Beyond the support,
$U^\sigma$ is monotone, and beyond $y=60$ the tail bound $W_\sigma\le3\log(y+1)-4\log y-ty+4\log(y+R)$ applies.

For a bent frame, $\sigma$ is a sum of uniform measures on panels of $\Gamma$ and on $[-4,-1]$. The admissibility
$W_\sigma\le\theta$ is proved on cells along each edge, the energy is a finite sum of explicit panel integrals, and the
atom penalty is bounded by a piecewise linear majorant of $\omega-2U^\sigma(-x)$, which covers every $\rho$ at once.

Every such inequality is a finite computation with rational numbers and rational enclosures of logarithms. These were
designed with the interval arithmetic of Arb \cite{Arb} and are checked in the formal proof by the Lean kernel. The
data were found by numerical optimization and play no other role in the proof.

%% =====================================================================
\section{Uniformity in the argument}\label{sec:uniform}

\subsection{Structure of the dual}
For a frame $F$ write $D^F_\rho(\sigma,\theta;t)$ for the dual value \eqref{eq:dual} with $t$ as a parameter, and call
the datum admissible at $t$ if $W_\sigma\le\theta$ on the contour.

\begin{lemma}\label{lem:U}
Fix a sign and a frame.
\begin{enumerate}[label=(\alph*)]
\item (concavity) The set of $(\sigma,\theta,t,\rho)$ with $(\sigma,\theta)$ admissible at $t$ is convex, and
 $D^F_\rho(\sigma,\theta;t)$ is jointly concave on it.
\item (transport) Let $0\le m\le1$, with $m=0$ for a bent frame. If $(\sigma,\theta)$ is admissible at $t$ and
 $\Delta\ge0$, then $(\sigma,\theta-m\Delta)$ is admissible at $t+\Delta$, and
 \[
  D^F_{\rho+(4+m)\Delta}(\sigma,\theta-m\Delta;t+\Delta)\ge D^F_\rho(\sigma,\theta;t)+(2m-m^2)\Delta .
 \]
\item (monotonicity) $D^F_\rho(\sigma,\theta;t)$ is nondecreasing in $\rho$.
\item (sign) For the lattice frame the dual value and admissibility do not depend on the sign.
\end{enumerate}
\end{lemma}

\begin{proof}
(a) The fields $W_\sigma-\theta$ and $A_\rho-\theta$ are affine in $(\sigma,\theta,t,\rho)$, so the positive parts and their
integrals are convex, and admissibility is a convex condition. For $\sigma_\lambda=\lambda\sigma_0+(1-\lambda)\sigma_1$,
\[
 I(\sigma_\lambda)=\lambda I(\sigma_0)+(1-\lambda)I(\sigma_1)-\lambda(1-\lambda)I(\sigma_0-\sigma_1),
\]
and $I(\sigma_0-\sigma_1)\ge0$ for a neutral signed measure of finite energy \cite[Lemma~I.1.8]{SaffTotik}. So $-I$ is
concave.

(b) At $t+\Delta$ the fields are $W_\sigma-\Delta\operatorname{Re}z$ and $A_\rho+\Delta x$. On a bent frame
$\operatorname{Re}z\ge0=m$. On the lattice,
\[
 \int_0^\infty\bigl(W_\sigma-\Delta y-\theta+m\Delta\bigr)_+dy\le\int_0^\infty(W_\sigma-\theta)_+dy+\Delta\int_0^\infty(m-y)_+dy ,
\]
and the last integral is $m^2/2$. With $\rho+(4+m)\Delta$ in place of $\rho$, the atom field changes by
$\Delta(x-4)\le0$ on $[1,4]$. The level term gains $2m\Delta$.

(c) $A_\rho$ decreases in $\rho$. (d) The lattice frame uses only $|\lambda|$.
\end{proof}

By Lemma~\ref{lem:U}(b), lowering the level by $\Delta$ gains two units of mass, and the saturated lattice loses at most
$\Delta\int_0^12(1-y)\,dy=\Delta$. This is the source of the slope $1$ in $t$.

\subsection{The table}
A \emph{node} is a frame together with a rational number $\lambda_k$ with $0<|\lambda_k|\le\frac12$. Its value
$t_k=\log(1/|\lambda_k|)$ is the parameter of its entries. An \emph{entry} of a node is a pair of rationals
$(\rho_k,D_k)$ for which a dual datum admissible at $t_k$ with $D^F_{\rho_k}\ge D_k$ is certified as in
Section~\ref{ssec:cert}. Our table has $24$ lattice nodes and $13$ bent nodes, with $151$ and $268$ entries. Lattice
nodes serve both signs, and bent nodes serve $s=-1$. Appendix~\ref{app:table} lists them.

\begin{definition}\label{def:Dlow}
For a sign $s$, $t>0$ and $\rho>0$, a number $D$ is a \emph{candidate} at $(s,t,\rho)$ if it arises in one of two ways
from entries of nodes serving the sign $s$.
\begin{enumerate}[label=(\Roman*)]
\item (interpolation) Two entries $(\rho_i,D_i)$ and $(\rho_j,D_j)$ of nodes on the same frame with $t_i<t_j$,
 $t_j-t_i\le\frac45$, $t=\theta t_i+(1-\theta)t_j$ with $\theta\in[0,1]$, and $\theta\rho_i+(1-\theta)\rho_j\le\rho$. Then
 $D=\theta D_i+(1-\theta)D_j$.
\item (transport) One entry $(\rho_k,D_k)$ with $t_k\le t$ and a slope $m\in\{0,\frac14,\frac12,\frac34,1\}$ if the node
 is a lattice node, or $m=0$, with $\rho_k+(4+m)(t-t_k)\le\rho$. Then $D=D_k+(2m-m^2)(t-t_k)$.
\end{enumerate}
Let $\Dlow_s(t,\rho)$ be the largest candidate, or $-\infty$ if there is none.
\end{definition}

\begin{theorem}[uniform decay]\label{thm:U}
Let $\lambda$ be rational with $0<|\lambda|<1$, $s=\operatorname{sign}\lambda$ and $t=\log(1/|\lambda|)$. If $D$ is a
candidate at $(s,t,\rho)$ with $\rho>0$, then for every $\varepsilon>0$ and all large $n$,
$|\tilde Q_n(\xi+u)|\le e^{-(D-\varepsilon)n^2}$ whenever $|u|\le e^{-(\rho+\varepsilon)n}$.
\end{theorem}

\begin{proof}
For an interpolation candidate, the mixture of the two certified data with weights $\theta$ and $1-\theta$ is
admissible at $t$ and has dual value at least $D$ at $\theta\rho_i+(1-\theta)\rho_j$, by Lemma~\ref{lem:U}(a). For a
transport candidate the transported datum of Lemma~\ref{lem:U}(b) does the same. Lemma~\ref{lem:U}(c) passes to $\rho$,
and Lemma~\ref{lem:U}(d) allows lattice nodes for $s=-1$. Theorem~\ref{thm:LD} with $\rho+\varepsilon$ in place of $\rho$,
and the monotonicity in $\rho$, give the claim.
\end{proof}

%% =====================================================================
\section{Proof of Theorem~\ref{thm:main}}\label{sec:proof}

\begin{proof}[Proof of Theorem~\ref{thm:main}]
Let $\lambda=a/b$ and $P_n=\tilde d_n\tilde Q_n\in\Z[X]$, of degree at most $2n$. Since $0<|\lambda|\le\frac12$, Theorem~\ref{thm:types}
shows that all type constants are nonzero at $\lambda$, and Corollary~\ref{cor:band} gives a band nonvanishing datum on
$(\frac1{13},\frac43)$ for $(P_n)$. So $M=2(\frac43-\frac1{13})=\frac{98}{39}$.

Let $D=\Dlow_s(t,\rho)$, attained by a candidate. By Theorem~\ref{thm:U}, $|\tilde Q_n(\xi+u)|\le e^{-(D-\varepsilon)n^2}$
for $|u|\le e^{-(\rho+\varepsilon)n}$ and $n$ large. By Theorem~\ref{thm:content}, $\tilde d_n\le e^{-(c(a,b)-\varepsilon)n^2}$
whenever $\tilde Q_n\ne0$. When $\tilde Q_n=0$, $P_n=0$. Hence
\[
 |P_n(\xi+u)|\le e^{-(\eta-2\varepsilon)n^2}\qquad(|u|\le e^{-(\rho+\varepsilon)n}),
\]
with $\eta=D+c(a,b)>0$. Theorem~\ref{thm:criterion} with $\kappa=2$ and $(\alpha,\beta)=(\frac1{13},\frac43)$ gives the
measure.
\end{proof}

\begin{proof}[Proof of Corollary~\ref{cor:table}]
Each entry is Theorem~\ref{thm:main} at one candidate. For example, at $\lambda=\frac13$ the lattice node $\frac13$ has
the entry $\rho=\frac{35}2$, $D=-2.624336\ldots$, and $c(1,3)=4.247432\ldots$. So $\eta=1.623095\ldots$ and
$\rho\max(2/\eta,39/98)=21.5637\ldots$. At $\lambda=-\frac13$ the bent node $-\frac13$ on $\Gamma_3$ has the entry
$\rho=18$, $D=0.336425\ldots$, and $c(-1,3)=2.687851\ldots$, which gives $11.9036\ldots$. At $\lambda=\frac12$ the entry
$\rho=\frac{31}2$, $D=-4.593612\ldots$ and $c(1,2)=4.725773\ldots$ give $\eta=0.132160\ldots$ and $234.5629\ldots$. At
$\lambda=-\frac2{39}$ the best candidate interpolates between the nodes $-\frac2{25}$ and $-\frac2{41}$ on $\Gamma_{15}$.
The other entries of Table~\ref{tab:values} are computed in the same way.

For Table~\ref{tab:larger-numerators}, let $u_i$ and $u_j$ be the magnitudes of the indicated nodes. For two
entries take $\theta=\log(|\lambda|/u_j)/\log(u_i/u_j)$ and interpolate both $D$ and $\rho$ with this weight.
The strict measure bounds follow by rational interval evaluation. Reduce each logarithm's argument to $[1,2]$ by
powers of two and use $\log x=2\sum_{k=0}^{31}v^{2k+1}/(2k+1)+R$, where $v=(x-1)/(x+1)$ and
$0\le R\le 2v^{65}/(65(1-v^2))$. For example entry $308$ has $\rho=91/4$ and
$D=5335835236283847/5000000000000000$ at $-2/11$. It gives $\eta>0.40891625$ and $\rho\max(2/\eta,39/98)<111.27$.

For the positive numerator $2$, transport entry $31$ from $2/13$ with slope $m=1$. The margin starts positive
at $b=13$, and its derivative in $b$ is $(17b-88)/(12b(b-2))>0$ for $b\ge13$. For numerator $-2$,
transport entry $308$ from $-2/11$ with slope $m=0$. The derivative of the margin is
$(5b+64)/(12b(b+2))>0$ for $b\ge11$. Thus the margins stay positive throughout the claimed ranges.
\end{proof}



%% =====================================================================
\section{Formalization}\label{sec:lean}
\raggedbottom

Theorem~\ref{thm:main} is formalized in Lean~4 \cite{Lean4} with Mathlib \cite{Mathlib}, together with every theorem,
proposition and lemma its proof uses. The remarks are not part of the formalization. The public statement is the
theorem \lean{Li2IrrationalityMeasure.li2_measure_general} in the statement file \lean{ChallengeLi2.lean} of the repository
\begin{center}
\FormalizationLink
\end{center}
Its statement uses only Mathlib. The number $\Li_2(a/b)$ is written as the series
$\sum_{k\ge0}(a/b)^{k+1}/(k+1)^2$, the rate $c(a,b)$ is written out, and the table of
Definition~\ref{def:Dlow} is given as rational literals. The conclusion is $q^{-\mu}<|\Li_2(a/b)-p/q|$ for all
integers $p$ and all sufficiently large $q$. The solution uses no \lean{sorry} and no \lean{native_decide}, and
\lean{\#print axioms} reports only \lean{propext}, \lean{Classical.choice} and \lean{Quot.sound}.

Tang's construction, his local bounds at the primes prime to $ab(b-a)$ and Proposition~\ref{prop:tang} are ported
from his formalization \cite{TangLi2} under the Apache License~2.0. The ported files are listed in the repository's
\texttt{NOTICE}. Tang's own files adapt parts of Firsching's formalization of Fauzan's proof \cite{Firsching26}. The
prime number theorem is taken from the PNT+ project \cite{PNTPlus}. The finite checks of Theorem~\ref{thm:types} and
of the certificates of Section~\ref{ssec:cert} are carried out by the Lean kernel in rational arithmetic. Their data
were computed with floating-point optimization and the interval arithmetic of Arb \cite{Arb}, and they play no role in
the proof beyond satisfying the checks.

\section*{Disclosure of AI use}
The mathematics, the Lean formalization, the certificate generators and the text of this paper were produced by AI
systems at the request of the author. The author posed the problems, directed the work and is responsible for the
paper. The systems were Claude Opus~5.5, Claude Sonnet~5.5 and Claude Haiku~5.5 (Anthropic), running in Claude Code,
and GPT-6.1 Sol (OpenAI), running in Codex. The construction and the irrationality proof of $\Li_2(1/2)$ are due to
Tang \cite{TangLi2}.

\subsection*{Acknowledgements}
We thank Qian Tang for releasing his formalization, the authors of Mathlib and of the PNT+ project, and Moritz
Firsching for his formalization of Fauzan's proof.

\appendix
\section{The table of nodes}\label{app:table}

Table~\ref{tab:nodes} summarizes the $37$ nodes of Definition~\ref{def:Dlow}: the node $\lambda_k$ ($\pm$ for a
lattice node, which serves both signs), its frame, $t_k$, the number of entries, the range of $\rho_k$ and the largest
$D_k$. The exact rational entries are part of the formal statement.

\begin{center}
\begin{longtable}{@{}lcrrcr@{}}
\caption{The nodes.}\label{tab:nodes}\\
\toprule
node & frame & $t_k$ & entries & $\rho_k$ & $\max D_k$\\
\midrule
\endfirsthead
\toprule
node & frame & $t_k$ & entries & $\rho_k$ & $\max D_k$\\
\midrule
\endhead
\bottomrule
\endfoot
$\pm1/2$ & lattice & 0.6931 & 6 & 14.5--16 & $-4.5932$ \\
$\pm1/3$ & lattice & 1.0986 & 1 & 17.5 & $-2.6243$ \\
$\pm1/4$ & lattice & 1.3863 & 1 & 18.5 & $-1.6650$ \\
$\pm1/5$ & lattice & 1.6094 & 1 & 20 & $-0.9524$ \\
$\pm1/6$ & lattice & 1.7918 & 1 & 20.5 & $-0.5379$ \\
$\pm2/13$ & lattice & 1.8718 & 14 & 22.3--23.55 & $-0.1476$ \\
$\pm2/15$ & lattice & 2.0149 & 7 & 22.5--24.5 & $0.2020$ \\
$\pm1/8$ & lattice & 2.0794 & 1 & 22 & $0.1472$ \\
$\pm2/17$ & lattice & 2.1401 & 6 & 23.5--25 & $0.4905$ \\
$\pm2/19$ & lattice & 2.2513 & 14 & 23.95--25.2 & $0.7348$ \\
$\pm1/10$ & lattice & 2.3026 & 1 & 22.9 & $0.5739$ \\
$\pm2/23$ & lattice & 2.4423 & 14 & 24.8--26.05 & $1.1294$ \\
$\pm2/25$ & lattice & 2.5257 & 15 & 25.2--26.45 & $1.2933$ \\
$\pm1/14$ & lattice & 2.6391 & 1 & 23.4 & $0.8978$ \\
$\pm2/33$ & lattice & 2.8034 & 14 & 26.45--27.7 & $1.8041$ \\
$\pm2/37$ & lattice & 2.9178 & 16 & 26.95--28.2 & $2.0011$ \\
$\pm2/39$ & lattice & 2.9704 & 18 & 27.2--28.45 & $2.0904$ \\
$\pm1/20$ & lattice & 2.9957 & 1 & 23.5 & $0.8591$ \\
$\pm1/26$ & lattice & 3.2581 & 3 & 24--30 & $2.5675$ \\
$\pm1/33$ & lattice & 3.4965 & 4 & 24.5--30.5 & $2.8818$ \\
$\pm1/55$ & lattice & 4.0073 & 3 & 28.5--33.5 & $3.6539$ \\
$\pm1/90$ & lattice & 4.4998 & 3 & 31--36 & $4.3096$ \\
$\pm1/148$ & lattice & 4.9972 & 3 & 33.5--38.5 & $4.9298$ \\
$\pm1/403$ & lattice & 5.9989 & 3 & 38.5--43.5 & $6.0988$ \\
$-1/2$ & $\Gamma_2$ & 0.6931 & 21 & 13.5--18.5 & $0.1226$ \\
$-1/3$ & $\Gamma_3$ & 1.0986 & 21 & 15--20 & $0.3465$ \\
$-1/4$ & $\Gamma_4$ & 1.3863 & 21 & 16--21 & $0.5109$ \\
$-1/5$ & $\Gamma_4$ & 1.6094 & 20 & 17.25--22 & $0.7657$ \\
$-2/11$ & $\Gamma_{15}$ & 1.7047 & 20 & 19.75--24.5 & $1.0672$ \\
$-1/6$ & $\Gamma_{15}$ & 1.7918 & 20 & 18.25--23 & $0.9669$ \\
$-2/15$ & $\Gamma_{15}$ & 2.0149 & 20 & 21--25.75 & $1.4102$ \\
$-2/25$ & $\Gamma_{15}$ & 2.5257 & 21 & 23--28 & $1.9449$ \\
$-1/13$ & $\Gamma_{15}$ & 2.5649 & 20 & 20--24.75 & $1.3212$ \\
$-1/20$ & $\Gamma_{15}$ & 2.9957 & 21 & 20--25 & $1.0338$ \\
$-2/41$ & $\Gamma_{15}$ & 3.0204 & 21 & 25--30 & $2.4113$ \\
$-2/65$ & $\Gamma_{15}$ & 3.4812 & 21 & 27.5--32.5 & $2.8658$ \\
$-1/33$ & $\Gamma_{15}$ & 3.4965 & 21 & 21.5--26.5 & $1.2118$ \\
\end{longtable}
\end{center}

\bibliographystyle{amsplain}
\bibliography{refs}

\end{document}
